\chapter{Deret Tak Hingga}

Sebuah tim melatih pengklasifikasi pada satu juta contoh. Fungsi objektifnya jumlah
kerugian tiap contoh, dan pelatihan berhenti otomatis ketika kenaikan total turun di
bawah $10^{-7}$ per epoch. Pada epoch keempat puluh proses berhenti sendiri. Modelnya
tetap buruk: setiap contoh baru hanya menyumbang sekitar $10^{-6}$ pada total, jadi
total hampir tidak bergerak, sementara total itu sendiri belum tentu menuju nilai
berhingga.

Kecurigaan itu bisa diuji dengan deret paling sederhana. Jumlah
$H_n = 1 + \tfrac{1}{2} + \cdots + \tfrac{1}{n}$ pada $n=10^6$ bernilai
$14.3927267229$. Menambah satu juta suku lagi hanya menaikkannya menjadi
$15.0858736534$; tambahannya $0.6931469306$, hampir tepat $\ln 2 = 0.6931471806$.
Tiap kali banyaknya suku berlipat dua, deret harmonik naik sekitar $\ln 2$, tanpa
batas: taksiran $\ln n + \gamma$ memberi $H_{10^9} \approx 21.3004815$. Aturan
berhenti berbasis "total sudah tidak bergerak" karena itu tidak membedakan konvergen
dari divergen; yang membedakan adalah laju pengecilan suku.

Bab 4 sudah memberi limit, kriteria Cauchy, kekonvergenan monoton, dan limit superior.
Deret tidak menambah gagasan baru, sebab jumlah parsial hanyalah barisan. Tetapi jumlah
parsial jarang punya rumus tertutup, dan menghitung $s_{10^6}$ tidak memberi tahu apakah
$s_{10^7}$ masih bergerak. Bab ini mengubah pertanyaan itu menjadi pertanyaan yang bisa
dijawab dengan memeriksa $a_n$ saja.

\begin{untukanak}
Bayangkan dua toples recehan. Di toples pertama kamu menabung dengan aturan aneh:
hari pertama satu koin, hari kedua setengah koin, hari ketiga seperempat koin, dan
seterusnya, selalu separuh dari hari sebelumnya. Kamu menabung setiap hari, tanpa
pernah libur, sampai kapan pun. Berapa isi toples itu setelah seratus tahun? Jawabannya
mengejutkan: tidak pernah lebih dari dua koin. Setiap tambahan memang mengecil, dan
jumlahnya berhenti di suatu angka yang bisa kamu sebut dengan pasti.

Di toples kedua aturannya berbeda. Koin hari pertama besar, lalu makin kecil, dan
makin kecil lagi. Setelah sebulan isinya baru segepok. Setelah setahun agak lebih
tebal. Kamu mungkin menyimpulkan toples ini juga akan berhenti di suatu angka, karena
tambahannya makin tipis. Ternyata tidak. Kalau kamu menabung cukup lama, isinya akan
melewati angka berapa pun yang kamu sebut, walaupun untuk naik satu angka saja kamu
butuh bertahun-tahun. Dua toples, dua aturan yang kelihatan mirip, dua nasib yang
berlawanan.

Pertanyaannya: tanpa menunggu bertahun-tahun, bagaimana tahu toples mana yang berhenti
dan mana yang tidak? Adakah ciri pada koinnya sendiri, bukan pada hasil akhirnya? Bab
ini menjawabnya, dan jawabannya dipakai setiap hari oleh orang yang menulis program.
\end{untukanak}

\section{Deret, Jumlah Parsial, dan Deret yang Bisa Dihitung Langsung}
\tingkat{dasar}

Deret\index{deret} adalah penjumlahan yang tidak pernah selesai. Karena itu
definisinya tidak memakai kata "jumlah" secara langsung, melainkan memakai barisan
jumlah sementara\index{jumlah parsial}.
Kita membaca hasil setelah sekian suku, lalu menanyakan ke mana bacaan itu menuju.

\begin{definisi}
Diberikan barisan $(a_n)$ di $\mathbb{R}$. Jumlah parsial ke-$n$ adalah
\[
  s_n = a_1 + a_2 + \cdots + a_n = \sum_{k=1}^{n} a_k .
\]
Lambang $\sum_{n=1}^{\infty} a_n$ disebut deret, dan deret itu \emph{konvergen} ke
$s$ bila barisan $(s_n)$ konvergen ke $s$; dalam hal ini kita menulis
$\sum_{n=1}^{\infty} a_n = s$ dan menyebut $s$ sebagai jumlah deret. Bila $(s_n)$
divergen, deret disebut \emph{divergen}.
\end{definisi}

\kamus{deret}{series}{penjumlahan yang tidak pernah selesai, yang nilainya dibaca dari jumlah sementara setelah sekian suku.}

Perhatikan $s_n - s_{n-1} = a_n$. Jadi deret dan barisan adalah benda yang sama; yang
berubah hanya pertanyaannya: pada barisan rumus $s_n$ diketahui, pada deret hanya $a_n$
yang kita pegang.

Kriteria Cauchy dari Bab 4 bisa langsung diterjemahkan, dan terjemahannya adalah alat
paling dasar di bab ini.

\begin{teorema}[Kriteria Cauchy untuk deret]
$\sum a_n$ konvergen jika dan hanya jika untuk setiap $\varepsilon > 0$ terdapat
bilangan bulat $N$ sehingga
\[
  m \ge n \ge N \quad \Longrightarrow \quad
  \left| \sum_{k=n}^{m} a_k \right| < \varepsilon .
\]
\end{teorema}

\noindent\textbf{Bukti.} Karena $\sum_{k=n}^{m} a_k = s_m - s_{n-1}$, syarat di atas tidak lain adalah syarat
Cauchy untuk barisan $(s_n)$. Menurut Bab 4, barisan real konvergen jika dan hanya
jika ia barisan Cauchy.

Kriteria ini bicara tentang potongan ekor deret, bukan tentang suku tunggal. Dengan
mengambil $m = n$ kita memperoleh syarat perlu yang sangat murah diperiksa.

\begin{akibat}[Uji suku ke-$n$]
Jika $\sum a_n$ konvergen, maka $a_n \to 0$.
\end{akibat}

\noindent\textbf{Bukti.} Ambil $m=n$ pada kriteria Cauchy: untuk setiap $\varepsilon>0$ terdapat $N$ sehingga
$n \ge N$ berakibat $|a_n| < \varepsilon$. Itu persis definisi $a_n \to 0$.

Uji ini hanya menyaring. Ia mengatakan bahwa suku yang tidak menuju nol pasti
menggagalkan konvergensi, dan tidak mengatakan apa pun tentang suku yang menuju nol.
Dua contoh cepat untuk sisi yang mudah: $\sum \frac{n}{n+1}$ divergen karena
$\frac{n}{n+1} \to 1$, dan $\sum (-1)^n$ divergen karena sukunya tidak menuju nol.
Sisi yang sulit adalah deret harmonik: $\frac{1}{n} \to 0$ tetapi $\sum \frac{1}{n}$
divergen. Kita buktikan itu di bagian kedua.

\begin{bacakeras}
Dua suara. Bacakan bergantian, lalu tukar perannya.

\textbf{Suara A:} Saya mengecil. Setiap hari lebih kecil daripada kemarin.

\textbf{Suara B:} Bagus. Tetapi mengecil belum berarti berhenti.

\textbf{A:} Saya mengecil separuh. Separuh dari separuh. Separuh lagi.

\textbf{B:} Maka kamu punya pagar. Jumlahmu akan berhenti.

\textbf{A:} Saya mengecil, tetapi hanya sedikit. Sepertiga. Seperempat. Seperlima.

\textbf{B:} Awas. Tipis bukan berarti habis.

\textbf{A:} Berapa pun kecilnya tambahan saya, saya tetap menambah.

\textbf{B:} Dan berapa pun lamanya, jumlahmu tetap naik.

Ulangi tiga kali makin cepat: mengecil itu urusan suku, berhenti itu urusan jumlah.
Hafalkan: yang menentukan bukan suku yang mengecil, melainkan sisa yang tidak pernah
habis.
\end{bacakeras}

Deret pertama yang bisa dihitung langsung adalah deret geometri\index{deret geometri}. Setiap sukunya
diperoleh dengan mengalikan suku sebelumnya dengan bilangan tetap.

\begin{teorema}[Deret geometri]
Untuk $|r| < 1$,
\[
  \sum_{k=0}^{\infty} r^k = \frac{1}{1-r},
\]
dan untuk $|r| \ge 1$ deret itu divergen.
\end{teorema}

\noindent\textbf{Bukti.} Untuk $r \neq 1$ berlaku identitas
\[
  (1-r)\left(1 + r + \cdots + r^n\right) = 1 - r^{n+1},
\]
sehingga $s_n = \frac{1-r^{n+1}}{1-r}$. Bila $|r|<1$ maka $r^{n+1} \to 0$ (Bab 4,
karena $|r|^{n} \to 0$ untuk $0 \le |r| < 1$), jadi $s_n \to \frac{1}{1-r}$. Bila
$|r| \ge 1$ maka $|r|^k \ge 1$ untuk semua $k$, sehingga $s_n$ tidak mungkin konvergen:
untuk $r=1$ jelas $s_n = n+1$, dan untuk $|r|\ge 1$ dengan $r\neq 1$ suku $r^k$ tidak
menuju nol.

Dengan $r = \tfrac{1}{2}$ kita memperoleh $1 + \tfrac{1}{2} + \tfrac{1}{4} + \cdots = 2$,
dan dengan $r = -\tfrac{1}{2}$ kita memperoleh $1 - \tfrac{1}{2} + \tfrac{1}{4} - \cdots
= \tfrac{2}{3}$. Yang kedua menarik: suku-sukunya berganti tanda dan jumlahnya lebih
kecil daripada suku pertama dikali dua, karena tanda yang berlawanan saling menghapus.

Kasus $r=\tfrac{1}{2}$ punya gambar yang menjelaskan dirinya sendiri: potong persegi
satuan menjadi dua, potong lagi sisanya, lalu teruskan. Gambar~\ref{fig:potong-geometri}
menunjukkan potongan $\tfrac12$, $\tfrac14$, $\tfrac18$ dan sisa $2^{-5} = 0.03125$
setelah lima potong, dengan luas terpakai $1 - 2^{-5} = 0.96875$. Sisa itu menyusut
tanpa lenyap; itulah rupa geometri dari kalimat "jumlahnya berhenti".

\begin{figure}[htbp]
\centering
\begin{tikzpicture}[scale=6.0]
  \draw[fill=layerteal!45,draw=ink] (0.5,0) rectangle (1,1);
  \draw[fill=layerindigo!40,draw=ink] (0,0.5) rectangle (0.5,1);
  \draw[fill=layerorange!50,draw=ink] (0.25,0) rectangle (0.5,0.5);
  \draw[fill=layerjade!45,draw=ink] (0,0.25) rectangle (0.25,0.5);
  \draw[fill=layerplum!40,draw=ink] (0.125,0) rectangle (0.25,0.25);
  \draw[fill=white,draw=ink,densely dashed] (0,0) rectangle (0.125,0.25);
  \draw[ink,thick] (0,0) rectangle (1,1);
  \node[font=\small] at (0.75,0.5) {$\tfrac{1}{2}$};
  \node[font=\small] at (0.25,0.75) {$\tfrac{1}{4}$};
  \node[font=\small] at (0.375,0.25) {$\tfrac{1}{8}$};
  \node[font=\small] at (0.125,0.375) {$\tfrac{1}{16}$};
  \draw[ink,thin,-{Stealth}] (0.06,0.10) -- (-0.30,-0.05);
  \node[font=\scriptsize,anchor=east] at (-0.30,-0.05) {sisa $\tfrac{1}{32}$};
\end{tikzpicture}
\caption{Deret geometri separuh sebagai potongan persegi satuan. Potongan berukuran
$\tfrac12$, $\tfrac14$, $\tfrac18$, $\tfrac1{16}$ dipasang berurutan, dan sisa yang
belum tertutup setelah empat potong adalah $\tfrac1{16}$; setelah lima potong sisanya
$2^{-5} = 0.03125$ dan luas terpakai $1 - 2^{-5} = 0.96875$. Sisa selalu menyusut,
tidak pernah lenyap.}
\label{fig:potong-geometri}
\end{figure}

\begin{tangan}
Bahan: satu lembar kertas persegi, pensil, penggaris.

\begin{enumerate}
\item Lipat kertas menjadi dua, buka, arsir separuhnya. Itu potongan pertama.
\item Lipat lagi bagian yang belum diarsir menjadi dua, buka, arsir separuhnya.
\item Ulangi tiga kali lagi, selalu hanya pada bagian yang belum diarsir.
\item Periksa bagian putih yang tersisa: setelah lima lipatan tinggal sepertiga puluh
      dua kertas semula.
\end{enumerate}

Yang harus kamu lihat: setiap arsiran baru selalu separuh dari yang sebelumnya, dan
bagian putih tidak pernah lenyap, hanya menguncup ke satu sudut. Kalau kertasnya cukup
tipis untuk dilipat sepuluh kali, bagian putih tinggal $\tfrac{1}{1024} = 0.0009765625$;
menurut hitungan ia menuju nol, sehingga arsiran akhirnya menutup seluruh kertas.
Sekarang ganti aturannya: arsir $\tfrac12$ kertas, lalu $\tfrac13$, lalu $\tfrac14$.
Lipatan tidak lagi bisa mengikuti, dan kertas habis sebelum deretnya berhenti.
\end{tangan}

Keluarga kedua yang bisa dihitung langsung adalah deret teleskopik\index{deret teleskopik}, yaitu deret yang
sukunya dapat ditulis sebagai selisih dua suku berurutan dari barisan lain sehingga
hampir semua pengurangan saling membatalkan.

\begin{contoh}
Untuk $n \ge 1$,
\[
  \frac{1}{n(n+1)} = \frac{1}{n} - \frac{1}{n+1}.
\]
Karena itu
\[
  s_N = \sum_{n=1}^{N} \left(\frac{1}{n} - \frac{1}{n+1}\right)
      = 1 - \frac{1}{N+1},
\]
sebab semua suku antara saling menghapus dan hanya suku pertama serta suku terakhir
yang tersisa. Jadi $\sum_{n=1}^{\infty} \frac{1}{n(n+1)} = 1$. Pada $N=100$ rumus itu
memberi $s_{100} = \tfrac{100}{101} = 0.9900990099$, dan sisa ekornya adalah
$\tfrac{1}{101} = 0.0099009901$; keduanya berjumlah tepat $1$.
\end{contoh}

Perhatikan bahwa pada contoh teleskopik kita tidak perlu uji apa pun. Rumus tertutup
untuk $s_N$ langsung memberi limit. Inilah yang jarang terjadi. Untuk hampir semua
deret, termasuk deret harmonik dan $\sum 1/n^2$, rumus seperti itu tidak ada, dan kita
harus memutuskan konvergensi tanpa pernah mengetahui jumlahnya secara tepat.

Sebelum masuk ke uji-uji, kita catat bahwa deret berperilaku linear, seperti halnya
limit barisan.

\begin{proposisi}
Jika $\sum a_n = A$ dan $\sum b_n = B$ konvergen, maka untuk setiap $c, d \in
\mathbb{R}$ deret $\sum (c a_n + d b_n)$ konvergen dengan jumlah $cA + dB$.
Menambahkan atau membuang berhingga banyak suku tidak mengubah konvergensi, hanya
mengubah jumlahnya.
\end{proposisi}

\noindent\textbf{Bukti.} Jumlah parsial deret gabungan adalah $c s_n^{(a)} + d s_n^{(b)}$, dan menurut sifat
aljabar limit di Bab 4 barisan itu konvergen ke $cA + dB$. Untuk pernyataan kedua,
misalkan suku dibuang mulai indeks $M$. Jumlah parsial deret baru adalah
$s_{N} - s_{M-1}$ untuk suatu $N$, yaitu selisih barisan konvergen dengan konstanta,
sehingga kekonvergenannya sama.

Sebagai pemakaian, $\sum_{n=1}^{\infty} \left(\frac{1}{2^n} +
\frac{1}{n(n+1)}\right) = 1 + 1 = 2$, karena bagian pertama adalah deret geometri
dengan $r=\tfrac{1}{2}$ yang dimulai dari $n=1$.

\section{Deret Suku Tak Negatif: Pembanding dan Deret-$p$}
\tingkat{menengah}

Bila semua suku tidak negatif, jumlah parsial hanya bisa naik atau diam. Pertanyaan
konvergen lalu berubah menjadi pertanyaan yang jauh lebih sederhana: apakah jumlah
parsial berhenti di bawah suatu pagar?

\begin{teorema}
Jika $a_n \ge 0$ untuk semua $n$, maka $\sum a_n$ konvergen jika dan hanya jika
barisan jumlah parsialnya terbatas.
\end{teorema}

\noindent\textbf{Bukti.} Dengan $a_n \ge 0$ berlaku $s_{n+1} = s_n + a_{n+1} \ge s_n$, jadi $(s_n)$ monoton
tak menurun. Menurut teorema kekonvergenan monoton di Bab 4, barisan monoton yang
terbatas di atas konvergen, dan barisan monoton yang tidak terbatas di atas divergen
menuju $+\infty$.

Teorema ini terlihat sederhana, tetapi seluruh bagian ini hanya variasi darinya. Yang
kita butuhkan hanyalah cara membandingkan deret yang belum dikenal dengan deret yang
sudah dikenal, dan persediaan deret yang sudah dikenal masih tipis: deret geometri.

\begin{teorema}[Uji pembanding]
Misalkan $0 \le a_n \le c_n$ untuk $n \ge N_0$.
\begin{enumerate}
\item Jika $\sum c_n$ konvergen, maka $\sum a_n$ konvergen.
\item Jika $a_n \ge d_n \ge 0$ untuk $n \ge N_0$ dan $\sum d_n$ divergen, maka
      $\sum a_n$ divergen.
\end{enumerate}
\end{teorema}

\noindent\textbf{Bukti.} Untuk (1), ambil $\varepsilon>0$. Menurut kriteria Cauchy ada $N \ge N_0$ sehingga
$m \ge n \ge N$ berakibat $\sum_{k=n}^{m} c_k < \varepsilon$. Karena setiap suku tidak
negatif dan $a_k \le c_k$,
\[
  \left| \sum_{k=n}^{m} a_k \right| = \sum_{k=n}^{m} a_k \le \sum_{k=n}^{m} c_k
  < \varepsilon ,
\]
jadi $\sum a_n$ memenuhi kriteria Cauchy. Untuk (2), bila $\sum a_n$ konvergen maka
dengan (1) dan $d_n \le a_n$ deret $\sum d_n$ ikut konvergen, bertentangan dengan
hipotesis.

Uji pembanding\index{uji pembanding} menuntut kita memilih pembanding yang tepat. Untuk suku berbentuk
pecahan rasional, pembanding yang wajar adalah suku terbesar di pembilang dibagi suku
terbesar di penyebut. Misalnya $\frac{1}{n^2 + 3n} \le \frac{1}{n^2}$, sehingga
$\sum \frac{1}{n^2+3n}$ konvergen asalkan $\sum \frac{1}{n^2}$ konvergen. Karena itu
seluruh pertanyaan berpindah ke satu keluarga deret.

Keluarga itu adalah deret-$p$\index{deret-p}, dan alat untuk menanganinya adalah pengelompokan suku
menjadi blok berukuran pangkat dua. Gagasan pengelompokan ini dipakai berulang kali,
termasuk pada uji kondensasi\index{kondensasi Cauchy} berikut.

\begin{lemma}[Kondensasi Cauchy]
Misalkan $a_1 \ge a_2 \ge a_3 \ge \cdots \ge 0$. Maka $\sum_{n=1}^{\infty} a_n$
konvergen jika dan hanya jika $\sum_{k=0}^{\infty} 2^k a_{2^k}$ konvergen.
\end{lemma}

\noindent\textbf{Bukti.} Karena sukunya tak negatif, kedua deret konvergen tepat ketika jumlah parsialnya
terbatas (teorema pertama bagian ini). Kita bandingkan blok-blok. Untuk $k \ge 0$,
blok $a_{2^k+1} + \cdots + a_{2^{k+1}}$ memuat $2^k$ suku, masing-masing tidak lebih
besar dari $a_{2^k}$ dan tidak lebih kecil dari $a_{2^{k+1}}$. Jadi
\[
  \sum_{n=1}^{2^m} a_n \;\le\; a_1 + \sum_{k=0}^{m-1} 2^k a_{2^k},
  \qquad
  \sum_{n=1}^{2^m} a_n \;\ge\; \frac{1}{2}\sum_{k=1}^{m} 2^k a_{2^k}.
\]
Pertidaksamaan kiri menunjukkan bahwa jumlah parsial $\sum a_n$ terbatas bila
$\sum 2^k a_{2^k}$ terbatas. Pertidaksamaan kanan menunjukkan sebaliknya, karena
jumlah parsial $\sum 2^k a_{2^k}$ terbatas bila $\sum a_n$ terbatas. Kedua arah
selesai.

\begin{teorema}[Deret-$p$]
$\sum_{n=1}^{\infty} \frac{1}{n^p}$ konvergen jika $p > 1$ dan divergen jika
$p \le 1$.
\end{teorema}

\noindent\textbf{Bukti.} Untuk $p \le 0$ suku $\frac{1}{n^p} = n^{-p}$ tidak menuju nol, jadi deret divergen
menurut uji suku ke-$n$. Untuk $p>0$ barisan $a_n = n^{-p}$ menurun menuju nol,
sehingga lemma kondensasi berlaku dan
\[
  \sum_{k=0}^{\infty} 2^k a_{2^k}
  = \sum_{k=0}^{\infty} 2^k \cdot 2^{-kp}
  = \sum_{k=0}^{\infty} \left(2^{1-p}\right)^{k}.
\]
Itu deret geometri dengan rasio $r = 2^{1-p}$. Karena $p>1$ tepat ketika $1-p<0$,
yaitu tepat ketika $r<1$, teorema deret geometri memberi: konvergen untuk $p>1$,
divergen untuk $p \le 1$.

Kasus $p=1$ adalah deret harmonik\index{deret harmonik}, dan buktinya sekarang hanya
satu baris: rasio
geometrinya $2^{1-1} = 1$, sehingga deret kondensasinya $\sum 1$ divergen. Kita bisa
membaca kerapuhannya langsung dari blok: setiap blok berukuran $2^k$ menyumbang paling
sedikit $2^k \cdot \frac{1}{2^{k+1}} = \tfrac{1}{2}$. Karena itu
\[
  H_{2^m} \;\ge\; 1 + \frac{m}{2},
\]
sehingga $H_n \to \infty$, tetapi untuk menaikkan $H$ sebesar satu perlu menggandakan
$n$ berkali-kali: laju pertumbuhannya logaritmik. Inilah alasan angka $14.39$ pada
pembukaan bab ini terlihat seperti batas padahal bukan.

Gambar~\ref{fig:batang-suku} menaruh kedua jenis suku berdampingan. Batang harmonik
tetap gemuk: suku kesepuluhnya masih $0.1$, sedangkan suku geometri kesepuluh
$2^{-10} = 0.0009765625$. Jumlah sepuluh suku pertama: $H_{10} = 2.928968254$ melawan
$1 - 2^{-10} = 0.9990234375$. Suku harmonik lebih besar pada setiap indeks, jadi besar
suku menyesatkan; yang menentukan adalah apakah pengecilannya meninggalkan ekor
berhingga.

\begin{figure}[htbp]
\centering
\begin{tikzpicture}
\begin{axis}[
  width=0.93\linewidth, height=6.0cm,
  ybar, bar width=3.4pt,
  xlabel={indeks $n$}, ylabel={besar suku $a_n$},
  xmin=0.3, xmax=10.7, ymin=0, ymax=1.12,
  xtick={1,2,3,4,5,6,7,8,9,10},
  scaled y ticks=false,
  tick label style={font=\scriptsize},
  label style={font=\small},
  ymajorgrids, grid style={layergrey!20},
  legend style={at={(0.98,0.96)},anchor=north east,draw=none,fill=none,
                font=\scriptsize},
]
\addplot[fill=layerteal!85,draw=layerteal] coordinates {
  (1,1) (2,0.5) (3,0.3333333) (4,0.25) (5,0.2)
  (6,0.1666667) (7,0.1428571) (8,0.125) (9,0.1111111) (10,0.1)};
\addlegendentry{harmonik $1/n$}
\addplot[fill=layerorange!85,draw=layerorange] coordinates {
  (1,0.5) (2,0.25) (3,0.125) (4,0.0625) (5,0.03125)
  (6,0.015625) (7,0.0078125) (8,0.0039063) (9,0.0019531) (10,0.0009766)};
\addlegendentry{geometri $2^{-n}$}
\end{axis}
\end{tikzpicture}
\caption{Batang suku harmonik $1/n$ dan geometri $2^{-n}$ untuk $n=1,\ldots,10$.
Batang harmonik lebih tinggi pada setiap indeks dan suku kesepuluhnya masih $0.1$,
sedangkan suku geometri kesepuluh $2^{-10} = 0.0009765625$. Jumlah sepuluh suku
pertama: $H_{10} = 2.928968254$ melawan $1 - 2^{-10} = 0.9990234375$.}
\label{fig:batang-suku}
\end{figure}

\kamus{deret harmonik}{harmonic series}{tumpukan yang terus bertambah tetapi tiap tambahannya makin tipis, sehingga tinggi totalnya naik sangat lambat tanpa batas.}

Deret-$p$ bukan batas akhir. Dengan kondensasi yang sama, Rudin menunjukkan bahwa
$\sum \frac{1}{n(\log n)^p}$ konvergen untuk $p>1$ dan divergen untuk $p<1$, dan bahwa
menambah satu logaritma lagi menggeser batasnya. Deret-deret itu berbeda sangat sedikit
tetapi berada di sisi berlawanan: tidak ada garis batas yang rapi, yang ada hanyalah uji
yang berlaku pada kelas tertentu.

\begin{tantangan}
Ambil $k(n)$ bilangan bulat dengan $2^{k-1} \le n < 2^k$, lalu tetapkan
$a_n = \frac{1}{k(n)}$. Jadi $a_1 = 1$, $a_2 = a_3 = \tfrac12$,
$a_4 = \cdots = a_7 = \tfrac13$, dan seterusnya. Tunjukkan bahwa $a_n \to 0$ tetapi
$\sum a_n$ divergen, lalu jelaskan mengapa uji rasio memberi limit superior $1$ di
sini sehingga tidak berguna.

\textbf{Petunjuk 1.} Jangan menjumlahkan suku satu per satu. Jumlahkan satu blok, dan
hitung dulu ada berapa suku di dalam blok itu.

\textbf{Petunjuk 2.} Blok ke-$k$ memuat $2^{k-1}$ suku yang semuanya bernilai
$\frac{1}{k}$, jadi sumbangannya $\frac{2^{k-1}}{k}$. Sepuluh blok pertama sudah
berjumlah $118.653968$.

Kalau mentok, kembali ke bagian Deret Suku Tak Negatif, tempat blok berukuran pangkat
dua dipakai pada Lemma kondensasi Cauchy.
\end{tantangan}

\begin{untukanak}
Coba bayangkan menumpuk koin di atas meja. Koin pertama tebal. Koin kedua setengah
tebalnya. Koin ketiga sepertiga, koin keempat seperempat, dan seterusnya: koin
ke-sejuta sangat tipis, setipis sepersejuta koin pertama. Kalau kamu menumpuk sejuta
koin, tinggi tumpukan itu kira-kira empat belas setebal koin pertama, bukan sejuta
kali. Kelihatannya tumpukan itu akan berhenti tumbuh, karena koin berikutnya hampir
tidak menambah apa pun.

Sekarang pertanyaan yang membuat banyak orang salah: berapa koin agar tumpukan
mencapai dua puluh kali tebal koin pertama? Jawabannya sekitar dua ratus tujuh puluh
dua juta koin. Untuk dua puluh satu kali, perlu sekitar sembilan miliar koin. Setiap
tingkat menuntut koin berkali-kali lipat, tetapi tidak ada tingkat yang mustahil.
Tumpukan itu tidak pernah berhenti, hanya melambat.

Inilah sebabnya program komputer mudah tertipu. Kalau kamu berhenti menumpuk setelah
sejuta koin, tumpukan tampak tidak bergerak lagi, padahal ia belum selesai. Di bab ini
kita akan tahu mengapa lambat bukan berarti berhenti, dan bagaimana membedakannya
sebelum menghabiskan waktu menumpuk.
\end{untukanak}

\section{Uji Rasio dan Uji Akar}
\tingkat{menengah}

Uji pembanding menuntut kita sudah punya deret pembanding. Untuk suku yang memuat
faktorial atau pangkat $n$, pembanding seperti itu sulit ditebak, tetapi perbandingan
antar suku berurutan mudah dihitung. Gagasannya: bila setiap suku paling banyak
$\beta$ kali suku sebelumnya dengan $\beta<1$, deretnya mengecil paling lambat seperti
deret geometri, jadi konvergen.

\begin{teorema}[Uji rasio]
Misalkan $a_n \neq 0$ untuk semua $n$ yang cukup besar.
\begin{enumerate}
\item Jika $\limsup_{n\to\infty} \left| \frac{a_{n+1}}{a_n} \right| < 1$, maka
      $\sum a_n$ konvergen, bahkan konvergen mutlak.
\item Jika terdapat $n_0$ sehingga $\left| \frac{a_{n+1}}{a_n} \right| \ge 1$ untuk
      semua $n \ge n_0$, maka $\sum a_n$ divergen.
\end{enumerate}
\end{teorema}

\noindent\textbf{Bukti.} Untuk (1), tulis $L = \limsup |a_{n+1}/a_n|$ dan pilih $\beta$ dengan $L < \beta < 1$.
Menurut sifat limit superior di Bab 4, hanya berhingga banyak $n$ yang memenuhi
$|a_{n+1}/a_n| > \beta$; jadi ada $N$ sehingga $n \ge N$ berakibat
$|a_{n+1}| \le \beta |a_n|$. Dengan induksi,
\[
  |a_{N+k}| \le \beta^{k} |a_N| \qquad (k \ge 0).
\]
Deret $\sum \beta^k |a_N|$ konvergen karena $0<\beta<1$, jadi menurut uji pembanding
$\sum |a_n|$ konvergen, dan menurut teorema konvergensi mutlak di bagian kelima bab ini,
kekonvergenan mutlak mengimplikasikan kekonvergenan.

Untuk (2), syarat $|a_{n+1}| \ge |a_n|$ untuk $n \ge n_0$ membuat $(|a_n|)$ tak
menurun dan tak negatif, sehingga tidak mungkin menuju nol. Maka $a_n$ tidak menuju
nol, dan uji suku ke-$n$ menggugurkan kekonvergenan.

\kamus{uji rasio}{ratio test}{mengukur apakah tiap suku menyusut paling sedikit sebesar sebagian tetap dari suku sebelumnya, sehingga deretnya mengecil seperti deret geometri.}
\index{uji rasio}

Batas $1$ pada uji rasio adalah titik buta. Untuk $1/n$ rasionya $\frac{n}{n+1} \to 1$
dan deretnya divergen; untuk $1/n^2$ rasionya $\left(\frac{n}{n+1}\right)^2 \to 1$
dan deretnya konvergen. Jadi rasio menuju $1$ tidak memberi informasi; orde pengecilan
yang berbeda jauh bisa memberi rasio yang sama.

Dua deret berikut memperlihatkan uji rasio bekerja pada arah yang berlawanan.
Untuk $a_n = \frac{n!}{10^n}$ berlaku
$\frac{a_{n+1}}{a_n} = \frac{n+1}{10}$, yang menuju $+\infty$; deretnya divergen.
Untuk $a_n = \frac{10^n}{n!}$ berlaku $\frac{a_{n+1}}{a_n} = \frac{10}{n+1}$, yang
menuju $0$; deretnya konvergen. Jumlah deret kedua tepat $e^{10} = 22026.4657948$,
sedangkan jumlah $20$ suku pertamanya baru $21950.3788494$.

Uji akar\index{uji akar} memakai besaran lain: akar ke-$n$ dari suku ke-$n$. Secara intuitif ia
mengukur "rasio geometri rata-rata" dari seluruh suku, bukan rasio antar dua suku
berurutan.

\begin{teorema}[Uji akar]
\begin{enumerate}
\item Jika $\limsup_{n\to\infty} \sqrt[n]{|a_n|} < 1$, maka $\sum a_n$ konvergen,
      bahkan konvergen mutlak.
\item Jika $\limsup_{n\to\infty} \sqrt[n]{|a_n|} > 1$, maka $\sum a_n$ divergen.
\end{enumerate}
\end{teorema}

\noindent\textbf{Bukti.} Untuk (1), pilih $\beta$ dengan $\limsup \sqrt[n]{|a_n|} < \beta < 1$. Hanya
berhingga banyak $n$ dengan $\sqrt[n]{|a_n|} > \beta$, jadi ada $N$ sehingga
$|a_n| \le \beta^n$ untuk $n \ge N$. Uji pembanding dengan deret geometri
$\sum \beta^n$ menutup kasus ini.

Untuk (2), bila limit superiornya lebih besar dari $1$, maka
$\sqrt[n]{|a_n|} > 1$, yaitu $|a_n| > 1$, untuk tak berhingga banyak $n$. Suku-suku
itu tidak menuju nol, jadi deretnya divergen menurut uji suku ke-$n$.

Uji akar lebih kuat daripada uji rasio. Misalkan $a_n = 2^{-n}$ untuk $n$ genap dan
$a_n = 3^{-n}$ untuk $n$ ganjil. Maka $\sqrt[n]{a_n}$ bernilai $\tfrac{1}{2}$ atau
$\tfrac{1}{3}$, sehingga limit superiornya $\tfrac{1}{2}$ dan uji akar menyatakan
deretnya konvergen. Rasionya berayun: pada $n=100$, $\frac{a_{101}}{a_{100}}$ bernilai
sekitar $3.05 \times 10^{17}$, dengan limit superior $+\infty$ dan limit inferior $0$.
Uji rasio tidak menyimpulkan apa pun, tepat ketika sukunya berayun dalam skala berbeda.

Uji rasio tetap populer karena faktorial dan pangkat $n$ saling membatalkan pada
rasio, sedangkan akar ke-$n$ dari faktorial menuntut aproksimasi Stirling. Urutan
pemakaian: uji suku ke-$n$, pembanding dengan deret-$p$ atau geometri, lalu uji rasio;
uji akar dipakai bila rasio berayun atau memberi $1$.

Sekarang uji-uji itu bertumpuk, dan urutan pemeriksaan menjadi keterampilan tersendiri.

\begin{resep}
\begin{enumerate}
\item Periksa $a_n \to 0$. Bila tidak, deret divergen; selesai.
\item Kenali bentuknya: faktorial memanggil uji rasio, pangkat $n$ memanggil uji akar,
      pecahan rasional memanggil pembanding deret-$p$.
\item Bila uji pilihanmu memberi batas $1$, jangan menyimpulkan apa pun. Pindah ke
      pembanding atau kondensasi.
\item Bila tanda berganti dan $|a_n|$ menurun ke nol, pakai Leibniz dan laporkan batas
      galat $b_{N+1}$.
\end{enumerate}

Peringatan urutan: langkah 1 jangan ditukar dengan langkah 2. Uji rasio pada deret yang
sukunya tidak menuju nol bisa memberi batas $1$ dan menyesatkan, sedangkan uji suku
ke-$n$ menyelesaikannya tanpa hitungan.
\end{resep}

\section{Deret Berganti Tanda dan Uji Leibniz}
\tingkat{menengah}

Sampai di sini semua deret bersuku tak negatif, dan di sana konvergensi berarti
jumlah parsial berhenti naik. Begitu tanda berganti-ganti, jumlah parsial bisa naik
dan turun, dan pembatalan antar tanda dapat menghasilkan konvergensi yang tidak
dijamin oleh besarnya suku.

\begin{teorema}[Uji Leibniz]
Misalkan $(b_n)$ barisan tak negatif yang tak menaik, yaitu $b_1 \ge b_2 \ge \cdots
\ge 0$, dan $b_n \to 0$. Maka $\sum_{n=1}^{\infty} (-1)^{n+1} b_n$ konvergen.
\end{teorema}

\noindent\textbf{Bukti.} Tulis $s_N$ untuk jumlah parsial ke-$N$. Pandang jumlah parsial berindeks genap dan
berindeks ganjil secara terpisah. Karena
\[
  s_{2m+2} - s_{2m} = b_{2m+1} - b_{2m+2} \ge 0,
\]
barisan $(s_{2m})$ tak menurun. Karena
\[
  s_{2m+1} - s_{2m-1} = -b_{2m} + b_{2m+1} \le 0,
\]
barisan $(s_{2m-1})$ tak menaik. Selain itu $s_{2m} \le s_{2m-1}$ dan
$s_{2m-1} - s_{2m} = b_{2m} \to 0$. Jadi $(s_{2m})$ tak menurun dan terbatas di atas
oleh $s_1 = b_1$, sehingga konvergen ke suatu $s$; $(s_{2m-1})$ tak menaik dan
terbatas di bawah oleh $s_2$, sehingga konvergen ke suatu $t$; dan
$s - t = \lim (s_{2m} - s_{2m-1}) = 0$. Kedua subbarisan menuju nilai yang sama, dan
karena subbarisan genap dan ganjil bersama-sama mencakup semua indeks, seluruh barisan
$(s_N)$ konvergen ke $s$ itu.

Uji Leibniz\index{uji Leibniz} memberi lebih dari konvergensi: batas galat yang bisa
dipakai tanpa mengetahui jumlah sebenarnya, dan itulah yang membuat deret berganti tanda
berguna untuk menghitung konstanta.

\begin{teorema}[Batas galat Leibniz]
Dengan hipotesis uji Leibniz dan $s = \sum (-1)^{n+1} b_n$,
\[
  \left| s - s_N \right| \le b_{N+1}.
\]
Lebih tajam: $s$ terletak di antara $s_N$ dan $s_{N+1}$, sehingga tanda galat
$s - s_N$ sama dengan tanda suku pertama yang dibuang, yaitu $(-1)^{N+2}$.
\end{teorema}

\noindent\textbf{Bukti.} Dari bukti sebelumnya, $(s_{2m})$ naik menuju $s$ dan $(s_{2m-1})$ turun menuju $s$.
Jadi untuk setiap $N$, bilangan $s$ berada di antara $s_N$ dan $s_{N+1}$. Karena
$|s_{N+1} - s_N| = b_{N+1}$, jarak $s$ ke $s_N$ tidak melebihi $b_{N+1}$, dan arah
pergeserannya ditentukan oleh tanda $s_{N+1} - s_N = (-1)^{N+2} b_{N+1}$.

Contoh klasiknya adalah deret harmonik berganti tanda,
\[
  1 - \frac{1}{2} + \frac{1}{3} - \frac{1}{4} + \cdots = \ln 2 .
\]
Kesamaan dengan $\ln 2$ dibuktikan lewat deret Taylor untuk $\ln(1+x)$ pada Bab 14;
di sini kita hanya memakai konvergensinya, yang dijamin uji Leibniz karena
$b_n = \frac{1}{n}$ menurun menuju nol. Angka $\ln 2$ muncul lagi, kini sebagai jumlah deret yang konvergen; selisih
$H_{2n} - H_n \to \ln 2$ adalah versi lain perhitungan yang sama.

Gambar~\ref{fig:tangga-leibniz} menggambar jumlah parsial deret itu sebagai dua tangga:
ganjil turun dari $s_1 = 1$, genap naik dari $s_2 = 0.5$, keduanya mengapit
$\ln 2 = 0.6931471806$. Pada $n = 10$ dan $n = 11$ apitannya
$0.6456349206 < \ln 2 < 0.7365440115$, selebar $b_{11} = \tfrac{1}{11} = 0.0909090909$,
tepat seperti yang dijanjikan teorema batas galat.

\begin{figure}[htbp]
\centering
\begin{tikzpicture}
\begin{axis}[
  width=0.93\linewidth, height=6.2cm,
  xlabel={banyak suku $n$}, ylabel={jumlah parsial $s_n$},
  xmin=0.4, xmax=15.0, ymin=0.45, ymax=1.06,
  xtick={1,2,3,4,5,6,7,8,9,10,11,12,13,14},
  tick label style={font=\scriptsize},
  label style={font=\small},
  ymajorgrids, grid style={layergrey!20},
  legend style={at={(0.98,0.04)},anchor=south east,draw=layergrey!40,fill=white,
                fill opacity=0.9,text opacity=1,font=\scriptsize},
]
\addplot[color=layerindigo,mark=*,mark size=1.6pt,thick] coordinates {
  (1,1.0) (3,0.8333333) (5,0.7833333) (7,0.7595238)
  (9,0.7456349) (11,0.7365440) (13,0.7301338)};
\addlegendentry{$n$ ganjil, turun}
\addplot[color=layerteal,mark=square*,mark size=1.6pt,thick] coordinates {
  (2,0.5) (4,0.5833333) (6,0.6166667) (8,0.6345238)
  (10,0.6456349) (12,0.6532107) (14,0.6587052)};
\addlegendentry{$n$ genap, naik}
\draw[dashed,accent,thick] (axis cs:0.4,0.6931472) -- (axis cs:15.0,0.6931472);
\node[font=\scriptsize,text=accent,fill=white,inner sep=1pt]
  at (axis cs:6.4,0.668) {$\ln 2 = 0.6931471806$};
\end{axis}
\end{tikzpicture}
\caption{Jumlah parsial deret harmonik berganti tanda. Subbarisan ganjil turun,
subbarisan genap naik, dan keduanya mengapit garis batas $\ln 2 = 0.6931471806$.
Pada $n = 10$ dan $n = 11$: $s_{10} = 0.6456349206$ dan $s_{11} = 0.7365440115$,
selebar $b_{11} = \tfrac{1}{11} = 0.0909090909$.}
\label{fig:tangga-leibniz}
\end{figure}

Batas galat Leibniz berguna dan sering ditaksir terlalu optimistis. Dengan $1000$
suku, jumlah parsialnya $0.7851481635$ sedangkan $\pi/4 = 0.7853981634$; galat
sebenarnya $2.5000 \times 10^{-4}$, batas yang dijanjikan
$b_{1001} = \tfrac{1}{2001} = 4.9975 \times 10^{-4}$. Batas itu benar tetapi dua kali
galat nyata, karena galat sebenarnya kira-kira $\tfrac{1}{4N}$ sedangkan batasnya
$\tfrac{1}{2N}$. Untuk galat di bawah $10^{-6}$ perlu $N \ge 500000$, dan pada $N$ itu
galatnya $5.0000 \times 10^{-7}$.

\section{Konvergensi Mutlak, Pengelompokan, dan Susunan Ulang}
\tingkat{lanjut}

Deret yang konvergen karena tandanya saling menghapus jauh lebih rapuh daripada deret
yang konvergen karena sukunya cukup cepat mengecil. Kerapuhan itu punya nama, dan
bab ini berutang satu penjelasan tentangnya.

\begin{definisi}
Deret $\sum a_n$ \emph{konvergen mutlak} bila $\sum |a_n|$ konvergen. Bila
$\sum a_n$ konvergen tetapi $\sum |a_n|$ divergen, deret itu disebut
\emph{konvergen bersyarat}\index{konvergensi bersyarat}.
\end{definisi}

\kamus{konvergensi mutlak}{absolute convergence}{deret yang tetap konvergen meskipun semua tandanya dipaksa menjadi positif, jadi tidak bergantung pada saling menghapus.}
\index{konvergensi mutlak}

\begin{teorema}
Jika $\sum a_n$ konvergen mutlak, maka $\sum a_n$ konvergen dan
$\left| \sum a_n \right| \le \sum |a_n|$.
\end{teorema}

\noindent\textbf{Bukti.} Untuk setiap $n \le m$ berlaku
$\left| \sum_{k=n}^{m} a_k \right| \le \sum_{k=n}^{m} |a_k|$. Karena $\sum |a_n|$
konvergen, ekspresi paling kanan bisa dibuat lebih kecil dari $\varepsilon$ menurut
kriteria Cauchy, sehingga $\sum a_n$ juga memenuhi kriteria Cauchy. Ketaksamaan
nilai mutlak jumlah diperoleh dengan mengambil limit pada
$\left| s_n \right| \le \sum_{k=1}^{n} |a_k|$.

Kebalikannya tidak berlaku: deret harmonik berganti tanda konvergen bersyarat,
karena $\sum \frac{1}{n}$ divergen. Pada deret seperti ini, urutan penjumlahan mulai
berperan. Sebelum sampai ke sana, kita lihat dulu operasi yang aman, yaitu
pengelompokan\index{pengelompokan}.

\begin{proposisi}[Pengelompokan]
Misalkan $\sum a_n$ konvergen ke $s$, dan $n_1 < n_2 < n_3 < \cdots$ barisan indeks.
Maka deret yang sukunya
\[
  A_k = a_{n_{k-1}+1} + a_{n_{k-1}+2} + \cdots + a_{n_k}
\]
(dengan $n_0 = 0$) juga konvergen ke $s$. Kebalikannya salah: pengelompokan yang
konvergen tidak menjamin deret asalnya konvergen.
\end{proposisi}

\noindent\textbf{Bukti.} Jumlah parsial deret kelompok adalah $S_K = A_1 + \cdots + A_K = s_{n_K}$, yaitu
subbarisan dari $(s_n)$. Subbarisan dari barisan konvergen menuju limit yang sama
(Bab 4), jadi $S_K \to s$.

Untuk kebalikannya, ambil $1 - 1 + 1 - 1 + \cdots$. Deret ini divergen karena sukunya
tidak menuju nol, tetapi dengan mengelompokkan berpasangan diperoleh
$(1-1) + (1-1) + \cdots = 0 + 0 + \cdots$, yang konvergen ke $0$.

Contoh penyangkal itu praktis penting: tanda kurung boleh dipasang pada deret yang sudah
diketahui konvergen, tetapi tidak boleh dipakai untuk membuktikannya.

Sekarang kita ubah urutannya, bukan hanya mengelompokkannya: susunan ulang menomori
kembali suku yang sama sehingga setiap suku muncul tepat sekali, tetapi pada posisi
berbeda.

\begin{definisi}
Misalkan $\sigma$ suatu bijeksi dari himpunan bilangan bulat positif ke dirinya
sendiri. Deret $\sum a_{\sigma(n)}$ disebut \emph{susunan ulang} dari $\sum a_n$.
\end{definisi}

\kamus{susunan ulang}{rearrangement}{menukar urutan kedatangan barang tanpa mengubah barangnya, sehingga total sementara bisa berbeda meski isi gudang sama.}
\index{susunan ulang}

Untuk deret yang konvergen mutlak, susunan ulang tidak berbahaya.

\begin{teorema}
Jika $\sum a_n$ konvergen mutlak, maka setiap susunan ulangnya konvergen, dan
semuanya menuju jumlah yang sama.
\end{teorema}

\noindent\textbf{Bukti.} Tulis $s = \sum a_n$ dan misalkan $\sum a_{\sigma(n)}$ susunan ulang dengan jumlah
parsial $s'_N$. Ambil $\varepsilon > 0$. Menurut kriteria Cauchy ada $N_0$ sehingga
$m \ge n \ge N_0$ berakibat $\sum_{k=n}^{m} |a_k| < \varepsilon$; khususnya
$\sum_{k=N_0}^{\infty} |a_k| \le \varepsilon$.

Pilih $P$ sehingga bilangan $1, 2, \ldots, N_0-1$ semuanya muncul di antara
$\sigma(1), \ldots, \sigma(P)$. Untuk $N > P$, suku-suku $a_1, \ldots, a_{N_0-1}$
saling menghapus pada selisih $s_N - s'_N$ (keduanya memuat suku itu), sehingga yang
tersisa hanyalah suku-suku berindeks $k \ge N_0$. Maka
\[
  \left| s_N - s'_N \right| \;\le\; \sum_{k \ge N_0} |a_k| \;\le\; \varepsilon
  \qquad (N > P).
\]
Karena $s_N \to s$ dan $\left| s'_N - s \right| \le \left| s'_N - s_N \right| +
\left| s_N - s \right|$, kita peroleh $s'_N \to s$.

Untuk deret yang hanya konvergen bersyarat, kesimpulan itu runtuh sepenuhnya. Teorema
berikut dirumuskan tanpa bukti; buktinya panjang dan teknis, dan dapat dibaca pada
Rudin, bab tentang barisan dan deret bilangan, bagian "Rearrangements".

\begin{teorema}[Susunan ulang Riemann, pernyataan]
Misalkan $\sum a_n$ konvergen tetapi tidak konvergen mutlak, dan misalkan $\alpha \le
\beta$ dua bilangan real (boleh $\pm\infty$). Maka ada susunan ulang $\sum
a_{\sigma(n)}$ yang jumlah parsialnya mempunyai limit inferior $\alpha$ dan limit
superior $\beta$. Khususnya, untuk setiap $c \in \mathbb{R}$ ada susunan ulang yang
konvergen ke $c$, dan ada susunan ulang yang divergen.
\end{teorema}

Gagasan buktinya begini. Bila $\sum a_n$ konvergen bersyarat, deret suku positifnya
dan deret nilai mutlak suku negatifnya sama-sama divergen. Kita menyusun ulang dengan
aturan: tambahkan suku positif sampai melewati $\beta$, lalu suku negatif sampai turun
di bawah $\alpha$, lalu ulangi. Karena kedua deret divergen, langkah itu selalu bisa
dilakukan, dan karena $a_n \to 0$, loncatan melewati batas semakin kecil. Hasilnya
osilasi dengan limit inferior $\alpha$ dan limit superior $\beta$. Konstruksi
lengkapnya ada pada Rudin, Teorema 3.54.

Kasus paling terkenal adalah deret harmonik berganti tanda yang disusun ulang dengan
pola dua suku positif lalu satu suku negatif,
\[
  1 + \frac{1}{3} - \frac{1}{2} + \frac{1}{5} + \frac{1}{7} - \frac{1}{4} + \cdots
\]
Deret ini konvergen, tetapi bukan ke $\ln 2$: jumlahnya $\tfrac{3}{2}\ln 2 =
1.0397207708$, dan perhitungannya ada di bagian berikut. Mengubah urutan mengubah
jawaban meskipun sukunya sama persis. Penjumlahan bilangan real tidak komutatif tanpa
syarat, dan itu peringatan konkret bagi siapa pun yang menjumlahkan pecahan di
komputer.

Ini juga menjelaskan mengapa uji di bagian kedua dan ketiga menuntut suku tak negatif.
Ketika tanda berganti, mengecil tidak cukup: pengecilannya harus cukup cepat mengalahkan
pembatalan.

\begin{untukanak}
Bayangkan dua tumpuk kartu: hijau bertanda "masuk", merah bertanda "keluar",
masing-masing dengan angka. Kamu memindahkan kartu satu per satu, dan saldo di layar
berubah: hijau menambah, merah mengurangi. Kalau urutan pemindahannya kamu ubah,
apakah saldo akhirnya berubah?

Kalau jumlah hijau dan merah sama-sama besar sehingga saling menghapus, urutannya
penting. Memindahkan hijau dulu membuat saldo melonjak sebelum turun; memindahkan
merah dulu membuat saldo sempat negatif. Saldo akhirnya pun bisa berbeda.

Ada satu keadaan aman: bila jumlah seluruh angka tanpa memandang warna tetap terbatas,
urutan tidak berpengaruh. Itulah konvergensi mutlak. Di bab ini kita akan tahu mengapa,
dan mengapa komputer bisa memberi jawaban berbeda untuk penjumlahan yang tampak sama.
\end{untukanak}

\section{Contoh Hitung}
\tingkat{menengah}

Bagian ini mengerjakan tiga hitungan yang menjadi tulang punggung bab. Semua angka
diperiksa dengan Python pada bagian lab; di sini langkah antaranya ditulis agar bisa
diulang dengan kertas dan pensil.

\medskip\noindent\textbf{1. Deret harmonik pada satu juta suku}

Kita ingin $H_{10^6}$ tanpa menjumlahkan satu juta suku. Alatnya adalah bentuk
asimtotik
\[
  H_n = \ln n + \gamma + \frac{1}{2n} - \frac{1}{12n^2} + O\left(\frac{1}{n^4}\right),
\]
dengan $\gamma = 0.5772156649$ adalah konstanta Euler--Mascheroni. Untuk $n = 10^6$:
\[
  \ln 10^6 = 13.8155105580, \qquad \ln 10^6 + \gamma = 14.3927262229 .
\]
Suku koreksi pertama $\frac{1}{2n} = 5 \times 10^{-7}$ memberi
$14.3927267229$, dan koreksi kedua $-\frac{1}{12n^2} = -8.33 \times 10^{-14}$
jauh di bawah presisi yang kita butuhkan. Jadi
\[
  H_{10^6} = 14.3927267228657\ldots
\]
Penjumlahan langsung dengan \texttt{math.fsum} memberi $14.392726722865724$, sepakat
sampai $10^{-15}$.

Yang lebih penting adalah berapa banyak suku yang dibutuhkan untuk mencapai batas
tertentu. Dari $H_n \approx \ln n + \gamma$ kita peroleh $n \approx e^{H - \gamma}$.
Untuk $H = 20$:
\[
  n \approx e^{20 - 0.5772156649} = e^{19.4227843351} = 2.724006 \times 10^{8}.
\]
Pencarian langsung memberi $n = 272400600$ sebagai indeks pertama dengan $H_n > 20$,
dan $H$ di titik itu $20.0000000016$. Untuk $H = 21$ perlu sekitar
$7.4 \times 10^8$ suku; taksiran pada $n = 10^9$ memberi $21.3004815$. Setiap
tambahan satu satuan menuntut penggandaan $n$ sebesar $e \approx 2.718$.

\medskip\noindent\textbf{2. Jumlah $\sum 1/n^2$ dan $\pi^2/6$}

Dengan $N = 10^6$ suku, penjumlahan langsung memberi
\[
  s_{10^6} = 1.6449330668487263, \qquad \frac{\pi^2}{6} = 1.6449340668482264 .
\]
Selisihnya $9.9999950 \times 10^{-7}$, sangat dekat dengan $\frac{1}{N} = 10^{-6}$;
itu masuk akal karena ekor deret $\sum_{n>N} n^{-2}$ kira-kira $\frac{1}{N}$.
Koreksi orde berikutnya adalah $-\frac{1}{2N^2} = -5 \times 10^{-13}$, sehingga
\[
  s_{10^6} + \frac{1}{N} - \frac{1}{2N^2} = 1.6449340668482262,
\]
yang cocok dengan $\pi^2/6$ sampai $2 \times 10^{-16}$, batas presisi pecahan
64-bit. Satu juta suku ditambah satu koreksi sudah memberi ketelitian mesin.

\medskip\noindent\textbf{3. Susunan ulang harmonik berganti tanda}

Ambil pola dua positif satu negatif: pada blok ke-$k$ kita menjumlahkan
$\frac{1}{4k-3} + \frac{1}{4k-1} - \frac{1}{2k}$. Setelah $K$ blok, seluruh suku
positif dengan penyebut $1, 3, \ldots, 4K-1$ sudah terpakai, dan suku negatif dengan
penyebut $2, 4, \ldots, 2K$ juga sudah terpakai. Jadi
\[
  S_K = \sum_{j=1}^{2K} \frac{1}{2j-1} - \sum_{j=1}^{K} \frac{1}{2j}
      = \left( H_{4K} - \frac{1}{2}H_{2K} \right) - \frac{1}{2}H_K .
\]
Dengan $H_m \approx \ln m + \gamma$,
\[
  S_K \approx \left(\ln 4K - \tfrac{1}{2}\ln 2K\right) - \tfrac{1}{2}\ln K
  = \ln \frac{4K}{\sqrt{2K}\sqrt{K}} = \ln\left(2\sqrt{2}\right) = \tfrac{3}{2}\ln 2 .
\]
Semua suku $\gamma$ menghapus, dan hasilnya tidak memuat $K$. Jadi
$S_K \to 1.0397207708$, bukan $\ln 2 = 0.6931471806$. Pada $K = 10^6$ blok nilai
numeriknya $1.0397205208$, galat $2.5 \times 10^{-7}$ terhadap $\tfrac{3}{2}\ln 2$,
sesuai orde $\frac{1}{4K}$ seperti deret Leibniz.

\kertas{Ringkasan angka. $H_{10^6} = 14.3927267229$ dan
$\ln 10^6 + \gamma = 14.3927262229$, selisihnya $\frac{1}{2n} = 5 \times 10^{-7}$.
$H_n > 20$ pertama kali tercapai pada $n = 272400600$, dekat
$e^{20-\gamma} = 2.724006 \times 10^8$. Deret Leibniz $1000$ suku memberi galat
$2.5000 \times 10^{-4}$ terhadap $\pi/4$, batasnya $\frac{1}{2001} =
4.9975 \times 10^{-4}$. Jumlah $\sum 1/n^2$ sampai $10^6$ suku kurang
$9.9999950 \times 10^{-7}$ dari $\pi^2/6$; satu koreksi Euler--Maclaurin menutupnya
sampai presisi mesin. Susunan ulang dua positif satu negatif menuju
$\frac{3}{2}\ln 2 = 1.0397207708$, bukan $\ln 2$.}

\section{Lab Komputasi}
\tingkat{lanjut}

Kode berikut dijalankan apa adanya pada Python 3.14.4 dengan NumPy 2.5.2. Bagian pertama
menguji ketelitian penjumlahan harmonik dan taksiran asimtotiknya.

\begin{lab}
\begin{verbatim}
import math
import numpy as np

n = 10**6
a = 1.0 / np.arange(1, n + 1)
vals = a.tolist()
naive = 0.0
for x in vals:
    naive += x

print("np.sum   ", repr(float(np.sum(a))))
print("for-loop ", repr(naive))
print("fsum     ", repr(math.fsum(vals)))
small = a.astype(np.float32)
print("float32  ", repr(float(np.sum(small, dtype=np.float32))))
g = 0.5772156649015328606
print("ln n + g ", repr(math.log(n) + g))
print("1/(2n)   ", 1 / (2 * n))
\end{verbatim}
\end{lab}

Keluaran yang diperoleh:
\begin{verbatim}
np.sum    14.392726722865726
for-loop  14.392726722864989
fsum      14.392726722865724
float32   14.392724990844727
ln n + g  14.392726222865807
1/(2n)    5e-07
\end{verbatim}

Empat hasil pertama seharusnya sama, dan tiga di antaranya berbeda pada digit
terakhir: galat pembulatan. Menjumlahkan sejuta bilangan berukuran $10^{-6}$ ke
akumulator berukuran $14$ membuang bit rendah. \texttt{float32} kehilangan
$1.7 \times 10^{-6}$, lebih besar daripada koreksi $\frac{1}{12n^2}$ yang ingin diukur;
bila galat yang dicari lebih kecil daripada galat penjumlahannya, hasilnya tidak
bermakna.

Bagian kedua memeriksa tiga angka dari bagian sebelumnya.

\begin{lab}
\begin{verbatim}
import math
import numpy as np

pi4 = math.pi / 4
s = 0.0
for k in range(1000):
    s += (-1)**k / (2*k + 1)
print("Leibniz 1000 suku", repr(s))
print("galat", repr(pi4 - s))
print("batas 1/2001", repr(1/2001))

N = 10**6
m = np.arange(1, N + 1, dtype=np.float64)
S = float(np.sum(1.0 / m**2))
print("jumlah 1/n^2", repr(S))
print("pi^2/6", repr(math.pi**2/6))
print("terkoreksi", repr(S + 1/N - 1/(2*N**2)))

S = 0.0
for k in range(1, 10**6 + 1):
    S += 1.0/(4*k-3) + 1.0/(4*k-1) - 1.0/(2*k)
print("susunan ulang", repr(S))
print("1.5 ln 2", repr(1.5*math.log(2)))
\end{verbatim}
\end{lab}

Keluarannya:
\begin{verbatim}
Leibniz 1000 suku 0.7851481634599485
galat 0.00024999993749974525
batas 1/2001 0.0004997501249375312
jumlah 1/n^2 1.6449330668487263
pi^2/6 1.6449340668482264
terkoreksi 1.6449340668482262
susunan ulang 1.0397205208399414
1.5 ln 2 1.0397207708399179
\end{verbatim}

Tiga hal perlu dibaca dari keluaran ini. Pertama, galat Leibniz $2.5000 \times
10^{-4}$ persis separuh batas $4.9975 \times 10^{-4}$; itu bukan kebetulan, karena
sisa deret berganti tanda yang monoton mendekati setengah suku pertama yang dibuang.
Kedua, koreksi Euler--Maclaurin menutup selisih sampai $2 \times 10^{-16}$, satu unit
pembulatan pada besaran $1.64$. Ketiga, susunan ulang memberi $1.0397205$, bukan
$0.6931472$, dan galatnya terhadap $\tfrac{3}{2}\ln 2$ hanya $2.5 \times 10^{-7}$.

\begin{untukdosen}
Urutan mengajar yang saya sarankan: mulai dari masalah epoch pada pembukaan, lalu
definisi jumlah parsial dan dua deret yang bisa dihitung (geometri, teleskopik).
Uji-uji masuk sebagai urutan pemeriksaan, bukan daftar hafalan: suku ke-$n$, pembanding,
rasio, akar. Deret-$p$ menyusul setelah kondensasi, yang menjelaskan mengapa harmonik
divergen dan mengapa lambat.

Dua miskonsepsi yang paling sering muncul. Pertama, mahasiswa menyimpulkan
konvergensi dari $a_n \to 0$; deret harmonik adalah obatnya, dan sebaiknya
ditunjukkan dengan tabel $H_{10} = 2.9290$, $H_{100} = 5.1874$, dan
$H_{1000} = 7.4855$ supaya lambatnya terlihat.
Kedua, mahasiswa menganggap "deret konvergen" berarti "jumlah parsial berhenti
berubah"; padahal perubahan per suku bisa jauh di bawah presisi mesin sementara
deretnya divergen. Ini tempat yang tepat untuk menyebut galat pembulatan.

Aktivitas kelas yang berhasil: bagi kelas menjadi kelompok, tiap kelompok diberi
empat deret tanpa nama, dan tugasnya hanya menentukan konvergen atau divergen
beserta ujinya dalam sepuluh menit. Deret yang dipilih harus memuat satu jebakan
rasio $=1$. Diskusikan kelompok yang memakai uji berbeda dan sama-sama benar.
Mahasiswa siap lanjut bila bisa menjelaskan mengapa $\sum 1/n^2$ konvergen sementara
$\sum 1/n$ tidak, cukup dengan pengelompokan blok.
\end{untukdosen}

\section{Ringkasan}
\tingkat{dasar}

\begin{ringkasan}
\begin{itemize}
\item Deret adalah barisan jumlah parsial; konvergen berarti barisan itu konvergen.
\item Untuk suku tak negatif, konvergen setara dengan jumlah parsial terbatas. Dari
      sini lahir uji pembanding, kondensasi, dan klasifikasi deret-$p$: konvergen
      tepat ketika $p>1$.
\item Uji rasio dan uji akar mengukur kecepatan pengecilan relatif terhadap deret
      geometri. Keduanya buta pada batas $1$, dan uji akar lebih kuat daripada uji
      rasio.
\item Uji Leibniz memberi konvergensi untuk deret berganti tanda yang menurun menuju
      nol, sekaligus batas galat $b_{N+1}$ yang tandanya diketahui.
\item Konvergensi mutlak kebal terhadap pengelompokan dan susunan ulang. Konvergensi
      bersyarat tidak: susunan ulang Riemann dapat mengubah jumlahnya menjadi
      bilangan apa pun, atau membuatnya divergen.
\item Deret harmonik divergen secara logaritmik: $H_{10^6} \approx 14.39$ dan perlu
      $272$ juta suku untuk mencapai $20$. Aturan berhenti berbasis "total tidak
      berubah lagi" tidak dapat dipercaya.
\end{itemize}
\end{ringkasan}

\section{Latihan}
\tingkat{menengah}

\begin{latihan}
\begin{enumerate}
\item Hitung $s_{20}$ untuk deret geometri $\sum_{k=0}^{\infty} \left(\frac{3}{5}\right)^k$,
      lalu bandingkan dengan jumlah sebenarnya. Berapa galatnya, dan berapa batas
      galat yang bisa kamu janjikan dari rumus $\frac{r^{21}}{1-r}$?

\item Buktikan bahwa $\sum_{n=1}^{\infty} \frac{1}{n(n+2)} = \frac{3}{4}$ dengan
      menguraikan suku menjadi pecahan parsial dan menunjukkan deretnya teleskopik.

\item Buktikan: bila $a_n \ge 0$ dan $\sum a_n$ konvergen, maka $\sum a_n^2$
      konvergen. Tunjukkan dengan contoh bahwa kebalikannya salah.

\item Nilai klaim ini: benar, salah, atau belum cukup bukti. "Karena
      $\frac{1}{n} \to 0$ dan setiap sukunya lebih kecil daripada suku deret
      geometri dengan rasio $\frac{1}{2}$ untuk $n$ cukup besar, deret harmonik
      konvergen." Sebutkan tepat di mana penalaran itu patah.

\item Tentukan konvergen atau divergen, dengan menyebut uji yang dipakai:
      (a) $\sum \frac{n^2}{2^n}$, (b) $\sum \frac{2^n}{n^2}$,
      (c) $\sum \frac{n!}{n^n}$, (d) $\sum \frac{1}{n^{1/3}}$.

\item Hitung jumlah $\sum_{n=1}^{10^6} \frac{1}{n^2}$ dengan NumPy, lalu taksir
      ekornya dengan $\frac{1}{N} - \frac{1}{2N^2}$ dan bandingkan hasilnya dengan
      $\pi^2/6$. Berapa angka di belakang koma yang benar sebelum koreksi, dan
      sesudahnya?

\item Deret $\sum \frac{(-1)^{n+1}}{\sqrt{n}}$ konvergen. Tentukan $N$ terkecil
      sehingga batas Leibniz menjamin galat di bawah $10^{-3}$, lalu hitung galat
      sebenarnya pada $N$ itu dengan Python. Apakah batasnya longgar atau ketat?

\item Buktikan bahwa bila $\sum a_n$ konvergen mutlak dan $(b_n)$ terbatas, maka
      $\sum a_n b_n$ konvergen mutlak.

\item Susun ulang deret harmonik berganti tanda dengan pola tiga positif satu
      negatif. Tebak jumlahnya dengan metode bagian contoh hitung, lalu periksa
      secara numerik. Apakah hasilnya masih berbentuk kelipatan $\ln 2$?

\item Buktikan bahwa bila $\sum a_n$ konvergen bersyarat, maka deret suku
      positifnya divergen dan deret nilai mutlak suku negatifnya juga divergen.
      Jelaskan mengapa fakta ini menjadi bahan bakar teorema susunan ulang Riemann.
\end{enumerate}

\noindent\emph{Petunjuk untuk tiga butir tersulit.} Butir 7: batas Leibniz memakai
suku pertama yang dibuang, jadi selesaikan $\frac{1}{\sqrt{N+1}} < 10^{-3}$ dan
periksa apakah ketaksamaan itu cukup tajam. Butir 9: setelah $K$ blok, suku positif
terpakai sampai penyebut $6K-1$, suku negatifnya sampai $\frac{1}{2K}$; tulis jumlah
parsialnya dalam $H_{6K}$, $H_{3K}$, dan $H_{K}$, lalu ambil limit. Butir 10: tulis
$p_n = \max(a_n, 0)$ dan $q_n = \max(-a_n, 0)$; maka $a_n = p_n - q_n$ dan
$|a_n| = p_n + q_n$, sehingga konvergensi salah satu deret itu menyimpulkan tentang
$\sum |a_n|$.
\end{latihan}

\begin{untukpeneliti}
Yang benar-benar terbukti di bab ini adalah pernyataan dengan bukti penuh: kriteria
Cauchy untuk deret, uji pembanding, kondensasi Cauchy, klasifikasi deret-$p$, uji rasio
dan uji akar dalam bentuk limit superior, uji Leibniz beserta batas galatnya, dan
ketertutupan deret konvergen mutlak terhadap susunan ulang. Pernyataan susunan ulang
Riemann saya kutip tanpa bukti. Sumber primernya Rudin, \emph{Principles of Mathematical
Analysis}, bab "Numerical Sequences and Series": Teorema 3.22--3.23, 3.25, 3.27--3.28,
3.33--3.34, 3.42--3.43, 3.45--3.46, dan 3.54--3.55. Abbott, \emph{Understanding
Analysis}, bab "Sequences and Series", membahas hal yang sama dengan penekanan pada
susunan ulang di bagian pembuka bab itu.

Yang sering dipopulerkan secara berlebihan adalah gagasan ambang batas antara deret
konvergen dan divergen. Rudin menunjukkan tidak ada ambang seperti itu: dua deret yang
berbeda sangat sedikit bisa berada di sisi berlawanan, dan prosedur itu bisa diteruskan
tanpa henti.

Pertanyaan terbuka: seberapa andal keputusan konvergensi bila $a_n$ hanya tersedia
sebagai taksiran berbising, seperti jumlah kerugian pada pembelajaran mesin? Untuk
deret semacam itu belum ada teori yang rapi.

Batas keberlakuan: seluruh bab membahas deret bilangan real. Di ruang bernorma lain
definisinya sama tetapi ujinya berbeda; deret fungsi menuntut kekonvergenan seragam
(Bab 14). Semua angka berasal dari NumPy 2.5.2, aritmetika 64-bit.
\end{untukpeneliti}

\peta{Bab 5 pada Rudin, \emph{Principles of Mathematical Analysis}, bab "Numerical
Sequences and Series": deret dan jumlah parsial (3.21--3.24), deret suku tak negatif
dan pembanding (3.25--3.29), uji akar dan rasio (3.33--3.37), konvergensi mutlak dan
susunan ulang (3.45--3.55). Pada Abbott, \emph{Understanding Analysis}, bab "Sequences
and Series": susunan ulang di bagian pembuka. Bab ini tidak memakai ulang prosa Bab 4;
dari sana hanya dirujuk definisi limit, kriteria Cauchy, dan kekonvergenan monoton.}
